%=============================================================================
% Section 4: Architectural Realization
%=============================================================================
\section{Architectural Realization}

\subsection{From Concept to Architecture}

While the contact layer is a conceptual model, it can be implemented through concrete system components. A practical architecture includes:

\begin{figure}[htbp]
\centering
\begin{tikzpicture}[
  layerbox/.style={
    draw, rectangle, minimum width=8cm, minimum height=0.8cm, align=center, font=\small
  },
  arrow/.style={
    -{Stealth[length=3mm, width=2mm]}, thick, gray!60!black
  }
]
% Layers
\node[layerbox, fill=blue!8] (agent) at (0, 4) {Agent (Intelligent Behavior)};
\draw[arrow] (0, 3.5) -- (0, 3);
\node[layerbox, fill=green!8] (contact) at (0, 2.5) {Contact Layer (Routing/Policy)};
\draw[arrow] (0, 2) -- (0, 1.5);
\node[layerbox, fill=orange!8] (exec) at (0, 1) {Execution Runtime (Execution)};
\draw[arrow] (0, 0.5) -- (0, 0);
\node[layerbox, fill=gray!8] (cap) at (0, -0.5) {Capabilities (External Actions)};
\end{tikzpicture}
\caption{Contact Layer Architecture Components}
\end{figure}

\textbf{Agent:} Represents intelligent behavior. Agents reason about tasks, plan actions, and decide which capabilities should be invoked.

\textbf{Contact Layer:} Performs capability routing and application policy decisions. This component determines which capability implementation to use, which provider should be selected, and which execution environment should handle the request.

\textbf{Execution Runtime:} Executes capability calls. This component handles the mechanics of invoking APIs, running tools, and returning results.

\textbf{Capabilities:} Represent actions available to the AI system.

\subsection{Clarifying Terminology Boundaries}

A common source of confusion is the relationship between Contact Layer and Execution Layer. Based on architectural clarity, we establish:

\textbf{Original Intuition:} Many practitioners conceive of a single ``execution/contact layer'' that handles both routing decisions and execution mechanics. This intuition is correct---there is one macro layer that mediates between agents and capabilities.

\textbf{Architectural Decomposition:} For software engineering purposes, this macro layer can be decomposed into two subcomponents:
\begin{enumerate}[label=\arabic*.]
\item \textbf{Contact Layer:} Handles routing, provider selection, and policy evaluation. This is the ``control plane'' of the action layer.
\item \textbf{Execution Runtime:} Handles API invocation, tool execution, and result normalization. This is the ``data plane'' of the action layer.
\end{enumerate}

\textbf{Unified Understanding:} The AI Action Layer (Macro Layer) comprises the Contact Layer (Routing/Policy - Control Plane) and the Execution Runtime (Execution - Data Plane). This decomposition maintains software engineering clarity while preserving the conceptual unity of the action interface.

\subsection{Precise Definition}

We can now state precisely:

\begin{quote}
\textit{The Contact Layer (in the narrow sense) is the routing and policy component of the AI Action Layer, responsible for capability routing, provider selection, and policy evaluation.}
\end{quote}

And more broadly:

\begin{quote}
\textit{The Action Layer (in the broad sense) is the complete interface between AI cognition and the capability ecosystem, comprising both routing (contact) and execution (runtime) components.}
\end{quote}

This terminology clarification resolves the apparent contradiction while preserving both conceptual clarity and architectural precision.
